\documentclass[9pt]{beamer}

\usecolortheme[RGB={205,173,0}]{structure} 
\usefonttheme{serif} 

\usepackage{amssymb}

%%%%%%%%%%%%%%%%%%%%%%%%%%%%%%%%%%%%%%%%%%%%%%%%
%%%%%%%%%%%%%%%%  Typesetting conventions  %%%%%%%%%%%%%%%%%%%

%doublebar
\def\double #1{#1{\hbox{\kern-2pt $#1$}}}
%Bras and kets
\def\ket#1{\left| #1\right>}
\def\bra#1{\left< #1\right|}
\def\braket#1#2{\left< #1\right|\left.#2\right>}
%For twistors
\def\kra#1{\left[ #1\right|}
\def\bet#1{\left| #1\right]}

%%%%%%%%%%%%%%%%%%%%%%%%%%%%%%%%%%%%%%%%%%%%%%%%
%%%%%%%%%%%%  Jim Gates and Marc Grisaru MACROS   %%%%%%%%%%%%%%%%
\def\pp{{\mathchoice
            %{general format
               %[w] = length of horizontal bars
               %[t] = thickness of the lines
               %[h] = length of the vertical line
               %[s] = spacing around the symbol
              %
              %\kern [s] pt%
              %\raise 1pt
              %\vbox{\hrule width [w] pt height [t] pt depth0pt
              %      \kern -([h]/3) pt
              %      \hbox{\kern ([w]-[t])/2 pt
              %            \vrule width [t] pt height [h] pt depth0pt
              %            }
              %      \kern -([h]/3) pt
              %      \hrule width [w] pt height [t] pt depth0pt}%
              %      \kern [s] pt
          {%displaystyle
              \kern 1pt%
              \raise 1pt
              \vbox{\hrule width5pt height0.4pt depth0pt
                    \kern -2pt
                    \hbox{\kern 2.3pt
                          \vrule width0.4pt height6pt depth0pt
                          }
                    \kern -2pt
                    \hrule width5pt height0.4pt depth0pt}%
                    \kern 1pt
           }
            {%textstyle
              \kern 1pt%
              \raise 1pt
              \vbox{\hrule width4.3pt height0.4pt depth0pt
                    \kern -1.8pt
                    \hbox{\kern 1.95pt
                          \vrule width0.4pt height5.4pt depth0pt
                          }
                    \kern -1.8pt
                    \hrule width4.3pt height0.4pt depth0pt}%
                    \kern 1pt
            }
            {%scriptstyle
              \kern 0.5pt%
              \raise 1pt
              \vbox{\hrule width4.0pt height0.3pt depth0pt
                    \kern -1.9pt  %[e]=0.15pt
                    \hbox{\kern 1.85pt
                          \vrule width0.3pt height5.7pt depth0pt
                          }
                    \kern -1.9pt
                    \hrule width4.0pt height0.3pt depth0pt}%
                    \kern 0.5pt
            }
            {%scriptscriptstyle
              \kern 0.5pt%
              \raise 1pt
              \vbox{\hrule width3.6pt height0.3pt depth0pt
                    \kern -1.5pt
                    \hbox{\kern 1.65pt
                          \vrule width0.3pt height4.5pt depth0pt
                          }
                    \kern -1.5pt
                    \hrule width3.6pt height0.3pt depth0pt}%
                    \kern 0.5pt%}
            }
        }}

\def\mm{{\mathchoice
                      %{general format %[w] = length of bars
                                       %[t] = thickness of bars
                                       %[g] = gap between bars
                                       %[s] = space around symbol
   %[w], [t], [s], [h]=3([g]) are taken from corresponding definitions of \pp
   %
                      %       \kern [s] pt
               %\raise 1pt    \vbox{\hrule width [w] pt height [t] pt depth0pt
               %                   \kern [g] pt
               %                   \hrule width [w] pt height[t] depth0pt}
               %              \kern [s] pt}
                  %
                       {%displaystyle
                             \kern 1pt
               \raise 1pt    \vbox{\hrule width5pt height0.4pt depth0pt
                                  \kern 2pt
                                  \hrule width5pt height0.4pt depth0pt}
                             \kern 1pt}
                       {%textstyle
                            \kern 1pt
               \raise 1pt \vbox{\hrule width4.3pt height0.4pt depth0pt
                                  \kern 1.8pt
                                  \hrule width4.3pt height0.4pt depth0pt}
                             \kern 1pt}
                       {%scriptstyle
                            \kern 0.5pt
               \raise 1pt
                            \vbox{\hrule width4.0pt height0.3pt depth0pt
                                  \kern 1.9pt
                                  \hrule width4.0pt height0.3pt depth0pt}
                            \kern 1pt}
                       {%scriptscriptstyle
                           \kern 0.5pt
             \raise 1pt  \vbox{\hrule width3.6pt height0.3pt depth0pt
                                  \kern 1.5pt
                                  \hrule width3.6pt height0.3pt depth0pt}
                           \kern 0.5pt}
                       }}

\def\ad{{\kern0.5pt
                   \alpha \kern-5.05pt
\raise5.8pt\hbox{$\textstyle.$}\kern 0.5pt}}

\def\bd{{\kern0.5pt
                   \beta \kern-5.05pt \raise5.8pt\hbox{$\textstyle.$}\kern 0.5pt}}

\def\qd{{\kern0.5pt
                   q \kern-5.05pt \raise5.8pt\hbox{$\textstyle.$}\kern 0.5pt}}
\def\Dot#1{{\kern0.5pt
     {#1} \kern-5.05pt \raise5.8pt\hbox{$\textstyle.$}\kern 0.5pt}}
%%%%%%%%%%%%%%%%%%%%%%%%%%%%%%%%%%%%%%%%%%%%%%%%%%%
%%%%%%%%%%%%%%%%%%%%%%%%%%%%%%%%%%%%%%%%%%%%%%%%%%%


\usepackage{beamerthemesplit}

\title{Covariant superstrings and compactification}
\author{William D. Linch, III}
\date{Overview of covariant string theory prepared for the High-energy Theory Group at Pontificia Universidad Cat\'olica de Chile.\\27 October, 2010}

\begin{document}

\frame{\titlepage}

\section[Outline]{}
\frame{\tableofcontents}

%%%%%%%%%%%%%%%%%%%%%%%%%%%%%%%%%%%%%%%%%%%%%%%%%%%
\section{Introduction}
\subsection{The role of string theory}
\frame{
\frametitle{Tools {\it versus} Physics}
The role of string theory is as a calculational tool for quantum field theories coupled to gravity.
\begin{eqnarray}
\begin{array}{ccc}
\textrm{TOOL}&\longleftrightarrow &\textrm{PHYSICS}\\
\textrm{Calculus}&&\textrm{classical mechanics}\\
\textrm{Functional Analysis}& &\textrm{quantum mechanics}\\
\textrm{QFT}& &\textrm{Fundamental fields with spin}\leq 1\\
\textrm{Strings}& &\textrm{Fundamental fields with spin}\leq2\\
\vdots & & \vdots
\end{array}
\nonumber
\end{eqnarray}
Note increasing symmetry and decreasing generality.
}
%%%%%%%%%%%
\subsection{How to use it?}
\frame
{
  \frametitle{Compactification and Holography}
\begin{itemize}
  \item Critical superstrings live in 10 flat dimensions
  \item We live in 4 (almost flat) dimensions
  \item To use our tool, we need to reduce $10\to 4$ 
  \item Two complementary possibilities:
  \item Compactification  
  \item Holography
  \end{itemize}
}

%%%%%%%%%%
\subsection{The need for covariant superstrings}
\frame{
\frametitle{Fluxes and branes}
Both approaches to dimensional reduction require analogue of electrons and/or background electromagnetic flux.

%\pause
{\bf Compactification}
%\pause
\begin{itemize}
 \item Compactifications on good string backgrounds without flux have moduli.
 %\pause
 \item Example: Heterotic string on Calabi-Yau 3-fold has $N=1$ supersymmetry $\Rightarrow$ flat directions.
\end{itemize}

%\pause
{\bf Holography}
%\pause
\begin{itemize}
 \item<1-> Holographic construction: Branes $\rightarrow$ Flux.
 %\pause
 \item<2-> Example: To get $N=4$ super-Yang-Mills we start with a stack of many D3-branes at weak string coupling. Increasing the coupling gives $AdS_5\times S^5$ with 5-form flux on the $S^5$.
\end{itemize}
}


%\subsection{Other applications}
%\frame{
%\frametitle{Current work}
%
%\begin{itemize}
% \item<1-> 
% \item<2-> 
% \item<3-> 
% \item<4-> 
%\end{itemize}
%}

%%%%%%%%%%%%%%%%%%%%%%%%%%%%%%%%%%%%%%%%%%%%%%%%%%%
\section{Superstrings}
\frame{
\frametitle{Superstrings}
Superstring = Bosonic string + SUSY
%\pause

Bosonic string $x:\Sigma \to M$
\begin{eqnarray}
S \sim \int ( \partial x)^2
\nonumber
\end{eqnarray}
%\pause
Supersymmetry in field theory
\begin{eqnarray}
\left.
\begin{array}{c}
\delta \phi  = \epsilon \psi  \\ 
\delta \psi =  \bar \epsilon\slash\!\!\! \partial \phi
\end{array}
\right\}
[\delta , \delta] \sim \bar \epsilon \gamma^m\epsilon \partial_m
\nonumber
\end{eqnarray}
}
\frame{
\frametitle{Superstrings}
Supersymmetry in string theory
\begin{eqnarray}
\begin{array}{cc}
\textrm {world-sheet} &\textrm {space-time} \\ 
\delta x^m = \epsilon \psi^m  &  \delta x^m = i\epsilon \gamma^m \theta \\ 
\delta \psi^m = \bar \epsilon \partial x^m & \delta \theta = \epsilon
\end{array}
\nonumber
\end{eqnarray}
%\pause
Open superstrings
\begin{eqnarray}
\begin{array}{cc}
\textrm {RNS$\pm$ spin structure:} &\textrm {GS$\pm$ spin structure:}\\
\textrm{(anti-)periodic B.C.} &\textrm {space-time (anti-)Weyl} \\ \\
S\sim \int (\partial x)^2 + \psi \partial \psi &  S\sim \int \Pi^2 +S_\mathrm{WZ} \\ 
			& \Pi^m = \mathrm d x^m +i \theta \gamma^m \mathrm d \theta
\end{array}
\nonumber
\end{eqnarray}
}
\frame{
\frametitle{Superstrings}
Closed superstrings
\begin{eqnarray}
\begin{array}{cccc}
	&\textrm {left-moving} &\textrm {right-moving}& \\ 
	& 1\pm &0& \textrm{Heterotic O/E}\\
\textrm{SUSY} &&\\	
	& 1\pm &1\pm &\textrm{Type II A/B}
\end{array}
\nonumber
\end{eqnarray}
}
\frame{
\frametitle{Superstrings}
Hidden symmetries
\begin{eqnarray}
\begin{array}{cc}
\textrm {RNS} &\textrm {GS} \\ \\
\textrm{Spacetime SUSY} & \textrm{Worldsheet Siegel symmetry} \\ \\
\left\{ \psi_0^m , \psi_0^n \right\} = 2g^{mn} & \delta \Pi^m = \mathrm d \bar \theta \gamma^m \delta \theta\\
	& \delta \theta = \slash\!\!\!\! \Pi \kappa \\ \\
\textrm{vacuum: spacetime spinor} & \textrm{generated by } \slash\!\!\!\! \Pi d \\ 
%\textrm{Bosonization } \bar \psi \psi \sim i\partial H & \textrm{Light-cone gauge} \\ 
%\left(\psi_0^m \pm J^m{}_n \psi_0^n\right) = \mathrm e^{i( \pm H_0\pm \dots \pm H_4)} & \gamma^+ \theta = 0 \\ 
%\downarrow \textrm{GSO-projection} & \Downarrow  \\ 
%\Theta_{\mathbf s} = \mathrm e^{i \sum_a s_a H_a} & x^i ~,~ S^a~~(\textrm{triality}) \\ 
%\textrm{Super-Virasoro} 	& \textrm{Unitary gauge} \\ 
%\psi \Theta \sim \frac1{\sqrt{z}} \dots 	& p_\theta = {\delta S\over \delta \Dot{ \theta}} \sim i \bar \theta \gamma^m \partial x_m+\dots
\end{array}
\nonumber
\end{eqnarray}
}
\frame{
\frametitle{Superstrings}
Hidden symmetries
\begin{eqnarray}
\begin{array}{cc}
\textrm {RNS} &\textrm {GS} \\ \\
%\textrm{Spacetime SUSY} & \textrm{Worldsheet Siegel symmetry} \\ 
%\left\{ \psi_0^m , \psi_0^n \right\} = 2g^{mn} & \delta \Pi^m = \mathrm d \bar \theta \gamma^m \delta \theta\\
%	& \delta \theta = \slash\!\!\!\! \Pi \kappa \\ 
%\textrm{vacuum: spacetime spinor} & \textrm{generated by } \slash\!\!\!\! \Pi d \\ 
\textrm{Bosonization } \bar \psi \psi \sim i\partial H & \textrm{Light-cone gauge} \\ 
\left(\psi_0^m \pm J^m{}_n \psi_0^n\right) = \mathrm e^{i( \pm H_0\pm \dots \pm H_4)} & \gamma^+ \theta = 0 \\ \\
\downarrow \textrm{GSO-projection} & \Downarrow  \\ \\
\Theta_{\mathbf s} = \mathrm e^{i \sum_a s_a H_a} & x^i ~,~ S^a~~(\textrm{triality}) \\ 
%\textrm{Super-Virasoro} 	& \textrm{Unitary gauge} \\ 
%\psi \Theta \sim \frac1{\sqrt{z}} \dots 	& p_\theta = {\delta S\over \delta \Dot{ \theta}} \sim i \bar \theta \gamma^m \partial x_m+\dots
\end{array}
\nonumber
\end{eqnarray}
}
\frame{
\frametitle{Superstrings}
Benefits/drawbacks
\begin{eqnarray}
\begin{array}{cc}
\textrm {RNS} &\textrm {GS} \\ \\
%\textrm{Spacetime SUSY} & \textrm{Worldsheet Siegel symmetry} \\ 
%\left\{ \psi_0^m , \psi_0^n \right\} = 2g^{mn} & \delta \Pi^m = \mathrm d \bar \theta \gamma^m \delta \theta\\
%	& \delta \theta = \slash\!\!\!\! \Pi \kappa \\ 
%\textrm{vacuum: spacetime spinor} & \textrm{generated by } \slash\!\!\!\! \Pi d \\ 
%\textrm{Bosonization } \bar \psi \psi \sim i\partial H & \textrm{Light-cone gauge} \\ 
%\left(\psi_0^m \pm J^m{}_n \psi_0^n\right) = \mathrm e^{i( \pm H_0\pm \dots \pm H_4)} & \gamma^+ \theta = 0 \\ \\
%\downarrow \textrm{GSO-projection} & \Downarrow  \\ \\
%\Theta_{\mathbf s} = \mathrm e^{i \sum_a s_a H_a} & x^i ~,~ S^a~~(\textrm{triality}) \\ 
\textrm{Super-Virasoro} 	& \textrm{Unitary gauge} \\ \\
\psi \Theta \sim \frac1{\sqrt{z}} \dots 	& p_\theta = {\delta S\over \delta {\dot \theta}} \sim i \bar \theta \gamma^m \partial x_m+\dots
\end{array}
\nonumber
\end{eqnarray}
}
\frame{
\frametitle{Superstrings}
Quantization 
\begin{eqnarray}
\begin{array}{cc}
\textrm {RNS} &\textrm {GS} \\ 
\textrm{Super-Virasoro constraints }1^\mathrm{st}\textrm{ class} & \textrm{Mixed }1^\mathrm{st}\textrm{ and }2^\mathrm{nd}\textrm{ class} \\ 
\textrm{Covariant quantization}  & \textrm{Non-covariant quantization}\\
\textrm{Trivial Ramond backgrounds}&\textrm{Trivial NS backgrounds}
\end{array}
\nonumber
\end{eqnarray}
{\em Neither possible for mixed backgrounds.}
}

%%%%%%%%%%%%%%%%%%%%%%%%%%%%%%%%%%%%%%%%%%%%%%%%%%%
\section{Covariant superstrings}
\frame{
\frametitle{Covariant superstrings}
%\subsection{Siegel strings} 
Origin of problem in covariant quantization of GS string: 
%\pause 

$2^\mathrm{nd}$ class constraints from composite $\theta$-momentum
%\pause

Siegel ``solution'': Include fundamental $\theta$-momentum
\begin{eqnarray}
\int (\partial x)^2  &\to &\int (\partial x)^2 + p_\theta \partial \theta
\nonumber
\end{eqnarray}
%\pause
Manifestly spacetime supersymmetric form
\begin{eqnarray}
S_\mathrm {Siegel} = S_\mathrm {GS}+ \int d \partial \theta 
\nonumber
\end{eqnarray}
%\pause
Siegel constraints ~~~~~$d =  p_\theta + i \bar \theta \gamma^m \partial x_m + \theta \theta \partial \theta\textrm{-term}$
}

\frame{
\frametitle{Covariant superstrings}
Operator products ~~~~~$d_\alpha (z)d_\beta(0) \sim \frac 1z (\gamma^m)_{\alpha \beta}\Pi_m(0)$
%\pause

~~~\\
Still mixed $1^\mathrm{st}$ and $2^\mathrm{nd}$ class since $T = \Pi^m \Pi_{m}+\dots$ is $1^\mathrm{st}$ class.
%\pause

~~~\\
Two solutions to the problem of separating constraints due to Berkovits:
\begin{enumerate}
%\pause
\item Hybrid strings
%\pause
\item Pure spinor strings
\end{enumerate}
}


%%%%%%%%%%%%%%%%%%%%%%%%%%%%%%%%%%%%%%%%%%%%%%%%%%%%%%%%%
\section{Hybrid strings}
\frame{
\frametitle{Hybrid strings}
Consider RNS in $M = \mathbb R^4 \times \mathbb C^3$ background
%\pause
\begin{eqnarray}
\begin{array}{lcc}
	&\textrm{``spacetime''}& \textrm{``internal''}\\\\
\textrm{matter}&x^m , \psi^m & y^i , \bar y^{\bar \imath}, \psi^i, \bar \psi^{\bar \imath}\\
\textrm{ghost} & b,c,\beta,\gamma &\\
%	&\downarrow&\downarrow\\
%\textrm{Berkovits variables}&x^m ,p_\alpha, \theta^\alpha,\bar p_{\dot \alpha}, \bar \theta^{\dot \alpha} ,\rho & y^i , \bar y^{\bar \imath}, \psi^i, \bar \psi^{\bar \imath}\\
\end{array}
\nonumber
\end{eqnarray}
%\pause
\begin{eqnarray}
\begin{array}{lcc}
%	&\textrm{``spacetime''}& \textrm{``internal''}\\\\
%\textrm{matter}&x^m , \psi^m & y^i , \bar y^{\bar \imath}, \psi^i, \bar \psi^{\bar \imath}\\
%\textrm{ghost} & b,c,\beta,\gamma &\\
\textrm{Berkovits}	&\downarrow&\downarrow\\
\textrm{variables}&x^m ,p_\alpha, \theta^\alpha,\bar p_{\dot \alpha}, \bar \theta^{\dot \alpha} ,\rho & y^i , \bar y^{\bar \imath}, \psi^i, \bar \psi^{\bar \imath}\\
\end{array}
\nonumber
\end{eqnarray}
%\pause
Examples
\begin{eqnarray}
(\theta^\alpha) &\sim& \exp \frac i2\left( -i \phi \pm (H_0+H_1)-H\right) \cr
(\bar \theta^{\dot \alpha})&\sim& c\xi \exp \frac i2 \left( 3 i \phi \pm (H_0-H_1) +H\right)
\nonumber
\end{eqnarray}
}
\frame{
\frametitle{Hybrid strings}
Canonical operator products
%\pause
\begin{eqnarray}
y^i(z,\bar z) \bar y^{\bar \jmath}(0) \sim \delta^{i\bar \jmath}\log|z|^2 
~,~ \psi^i(z) \bar \psi^{\bar \jmath}(0) \sim {\delta^{i \bar \jmath}\over z}
~,~ \rho(z)\rho(0) \sim -\log z
\cr
x^m(z,\bar z) x^n(0)\sim \log |z|^2
~,~p_\alpha (z) \theta^\beta(0) \sim {\delta_\alpha^\beta\over z} 
~,~\bar p_{\dot \alpha} (z) \bar \theta^{\dot \beta} (0) \sim {\delta_{\dot \alpha}^{\dot \beta}\over z}
\nonumber
\end{eqnarray}
%\pause
Super-Virasoro constraints
%\pause
\begin{eqnarray}
T& =& \Pi^m \Pi_m + d\partial \theta +\bar d \partial \bar \theta 
	+\delta_{i\bar \jmath} \partial y^i \partial \bar y^{\bar \jmath}
	+ \delta_{i\bar \jmath} \psi^i \stackrel{\leftrightarrow}{\partial} \bar \psi^{\bar \jmath}\cr
G^+&=& \mathrm e^{\rho} d^2 + i \delta_{i \bar \jmath}  \psi^i \partial \bar y^{\bar \jmath}\cr
G^{-}&=& \mathrm e^{-\rho} \bar d^2 + i \delta_{i \bar \jmath} \bar \psi^{\bar \jmath} \partial y^i\cr
J&=&-\partial \rho -\delta_{i \bar \jmath} \bar \psi^{\bar \jmath} \psi^i
\nonumber
\end{eqnarray}
}
%%%%%%%%%%%%
\frame{
\frametitle{Hybrid strings}
All constraints first class with commuting factors
%\pause
\begin{eqnarray}
\textrm{SCA}_{10} = \textrm{SCA}_4 \oplus \textrm{SCA}_6 &\Rightarrow& \textrm{Rep}_{10} = \textrm{Rep}_{4}\otimes \textrm{Rep}_{6}
\nonumber
\end{eqnarray}
%\pause
\begin{eqnarray}
\textrm{Rep}_{4} &\Leftrightarrow& \Phi (x, \theta, \bar \theta, p, \bar p, \rho)\cr\cr
\textrm{Rep}_{6} &\Leftrightarrow& \sum \Omega_{i_1 \dots i_p \bar \jmath_1\dots \bar \jmath_q}(y,\bar y) \psi^{i_1}\dots\psi^{i_p} \bar \psi^{\bar \jmath_1} \dots \bar \psi^{\bar \jmath_q}
\nonumber
\end{eqnarray}
}

\frame{
\frametitle{Physical state conditions}
Physical state conditions for unintegrated vertex operators are the usual RNS conditions. 
%\pause
%: 
%\begin{enumerate}
%\item No double poles with Virasoro 
%\item No single pole with $J$
%\end{enumerate}
No double poles with Virasoro implies conformal weight 0. 
%\pause
First possibility: $\Omega$ has weight 0 
%\pause
$\Rightarrow$ $\Omega =1$ 
%\pause
$\Rightarrow$ $\Phi$ function of 0-modes of $x,\theta, \bar \theta$ only 
%\pause
$\Rightarrow$ $\Phi$ is a {\em superfield} and 
%\pause
\begin{eqnarray}
d (z) \Phi  \sim \frac1z D \Phi &,& \bar d (z) \Phi  \sim \frac1z \bar D \Phi\cr\cr
D = {\partial\over \partial \theta} + i (\bar \theta \sigma^m){\partial \over \partial x^m}&,&
\bar D = -{\partial\over \partial \theta} - i (\sigma^m\theta){\partial \over \partial x^m}
\nonumber
\end{eqnarray}
Superspace covariant derivatives
%\pause
with algebra
\begin{eqnarray}
%d (z) \Phi  = \frac1z D \Phi && \bar d (z) \Phi  = \frac1z \bar D \Phi\cr\cr
\left\{D , D\right\}=0 ~,~\left\{D , \bar D\right\}=-2i \sigma^m \partial_m~,~ \left\{\bar D, \bar D\right\}=0
\nonumber
\end{eqnarray}
}

\frame{
\frametitle{(Anti-)chiral fields}
(Anti-)chiral fields: 

~~~\\
%\pause
No single pole with ($G^+$) $G^-$:
%\pause

~~~\\
$0 = \oint G^{-}U$
%\pause
$
= \left(\oint \mathrm e^{-\rho} \bar d^2(z) \Phi(0)\right)\Omega(0)+ \Phi(0) \left( \oint \bar \psi \partial y(z) \Omega(0)\right)$
%\pause
$
= \left(\mathrm e^{-\rho} \bar d\bar D\Phi\right)\Omega+\Phi\left( \bar \psi {\partial \Omega\over \partial \bar y} +\partial y {\delta \Omega\over \delta  \psi}\right)
$
%\pause

~~~\\
$\Rightarrow$ $\Phi$ is a chiral superfield and $\Omega$ is independent of $\psi$.

%\pause
$\Rightarrow \Omega = \omega(y, \bar y)_{\bar \imath_1\dots \bar \imath_q} \bar \psi^{\bar \imath_1}\dots \bar \psi^{\bar \imath_q}$ 
and 
$\bar \partial \Omega=0$ 

%\pause 
$\Rightarrow$ $\Omega$ holomorphic $(0,q)$-form!
%\pause
 
 ~~~\\
{\em ``Chiral ring'' of ``internal'' space $\Leftrightarrow$ cohomology ring}.
}

\frame{
\frametitle{(Anti-)chiral fields}
Spin-$\frac32$ generators 
%\pause
\begin{eqnarray}
\begin{array}{lllcl}
\textrm{Correlated}&G^+ = G^+_4 +G^+_6,&G^- = G^-_4 +G^-_6 &\textrm{and}& J = J_4+J_6\\
\textrm{or}&&&& \\
\textrm{Anti-corr.}&G^+ = G^+_4 +G^-_6,&G^- = G^-_4 +G^+_6 &\textrm{and}& J = J_4-J_6
\end{array}
\nonumber 
\end{eqnarray}
%\pause
Spacetime chiral superfields correlated with internal chiral ring 

{\em or} 

spacetime anti-chiral superfield correlated with chiral ring and 

{\it vice versa}.
}


%%%%%%%%%%%%%%
\subsection{Flux deformed Calabi-Yau compactification}
\frame{
\frametitle{Calabi-Yau compactification}

Type IIB string on Calabi-Yau $X$ 

%\pause
Left- and right-moving SCAs (Type II) 
%\pause

Both correlated (Type B)
%\pause
\begin{eqnarray}
\begin{array}{ccc}
	&(\textrm{chiral,chiral}) & (\textrm{chiral,anti-chiral})\\\\
\textrm{spacetime} & N=2\textrm{~chiral} &	N=2\textrm{~twisted-chiral}\\\\
\textrm{multiplet} & \textrm{vector} &\textrm{hyper}\\\\
\textrm{cohomology} & H^{2,1}_{\bar \partial}(X) & H^{1,1}_{\bar \partial}(X)\\\\
\textrm{deformation} &\mathbb C\textrm{-structure} &	\textrm{K\"ahler-structure}
\end{array}
\nonumber 
\end{eqnarray}
}

\frame{
\frametitle{Calabi-Yau compactification}

{\bf Vector-multiplets} (chiral,chiral) 
%\pause
$M = \sum_{p=1}^{h^{2,1}} M^p \Omega^p_{(-1,-1)}$
%\pause

\begin{eqnarray}
M^p = \phi^p + \theta_L^2 D^p_{++} + \theta_R^2 D_{--}^p + \theta_L\theta_R D_{+-}^p +\theta_L \sigma^{mn} \theta_R F^p_{mn}%+\theta_L^2\theta_R^2 \Box \bar \phi^p
+\dots 
\nonumber
\end{eqnarray}
%\pause

Integrated vertex operator 
%\pause

\begin{eqnarray}
I^p = \int G_L^+(G_R^+(M^p))+\mathrm {h.c.}= \int \mathrm e^{\rho_L +\rho_R} (d_{L} d_{R})D_{+-}^p+\mathrm {h.c.}+\dots
\nonumber
\end{eqnarray}
}

\frame{
\frametitle{Flux deformation}
D3-brane background preserves $\theta^{+} = \theta_L+i \theta_R$ spacetime supersymmetry
%\pause
\begin{eqnarray}
M^p = \phi^p + (\theta^+)^2 X^p_{++} + (\theta^-)^2 X_{--}^p + (\theta^+\theta^-) X_{+-}^p +\dots 
\nonumber
\end{eqnarray}
%\pause
Flux deformed background preserving D3 symmetry (Klebanov-Strassler)
%\pause
\begin{eqnarray}
X^p_{--} =0=X_{+-}^p &&  X_{++}^p \sim g_s \bar G^p ~~~(\textrm{ISD flux})
\nonumber
\end{eqnarray}
}

\frame{
\frametitle{Flux deformation}

Vertex operator for flux deformation 
%\pause
\begin{eqnarray}
I_{\mathrm D3}=\sum_{p=1}^{h^{2,1}} I^p\Omega^p = \int G_L^+(G_R^+(M^p))+\mathrm {h.c.}\cr
\sim g_s \int \mathrm e^{\rho_L +\rho_R} (d_{L} d_{R})\bar G^p\Omega^p +\mathrm{h.c.}
\nonumber
\end{eqnarray}
%\pause
%\begin{eqnarray}
%%I_{D3}=\sum_{p=1}^{h^{2,1}} I^p\Omega^p = \int G_L^+(G_R^+(M^p))+\mathrm {h.c.}
%\sim g_s \int \mathrm e^{\rho_L +\rho_R} (d_{L} d_{R})\bar G^p\Omega^p +\mathrm{h.c.}
%\nonumber
%\end{eqnarray}
%%\pause
Conformal symmetry breaking divergence
%\pause
\begin{eqnarray}
I_{\mathrm D3}I_{\mathrm D3} \sim g_s^2  \log\Lambda  \int \Pi_L\Pi_R|G|^2
\nonumber
\end{eqnarray}
%\pause
Cancel with {\em spacetime} gravitational vertex operator. 
}

\frame{
\frametitle{Warping}
In flat space
%\pause
\begin{eqnarray}
U_1 \sim (\theta_L\theta_R)(\bar\theta_L \bar \theta_R) A+\dots &\Rightarrow& I_1 = \int \Pi_L\Pi_R A+\dots
\nonumber
\end{eqnarray}

%\pause
Background field expansion 
%\pause
\begin{eqnarray}
I_1\to \int \Pi_L\Pi_R y^i \bar y^{\bar \jmath} \nabla_i \bar \nabla_{\bar \jmath} A \sim \log \Lambda \int \Pi_L\Pi_R \Delta A
\nonumber
\end{eqnarray}
%\pause
At 1-loop therefore 
%\pause 
~~~~~~~~~\fbox{$\Delta A \sim g_s^2 |G|^2$}

\vspace{.3cm}
{\em Warp factor of the spacetime metric} (\cite{Linch:2008nt} with Brenno Carlini Vallilo).
}

\frame{
\frametitle{Significance}
The results of the work with Brenno Carlini Vallilo have many implications:
\begin{itemize}
\item This approach allows a worldsheet ({\it i.e.} string-theoretic) derivation of many spacetime effects of background fluxes
%\pause
\item Can in principle be used to find all effects such as torsion, non-commutativity, {\it et cetera}
%\pause
\item Suggests the validity of the flux background as a string background due to existence of worldsheet SCA (no breaking of worldsheet supersymmetry)
%\pause
\item Allows computation of scattering amplitudes in non-trivial backgrounds
%\pause
\item Existence of new type of mirror symmetry for heterotic strings?
\end{itemize}
%\pause
Unfortunately, ``picture-changing problem'' persists... 
}

%%%%%%%%%%%%%%%%%%%%%%%%%%%%%%%%%%%%%%%%%%%%%%%%%%%%%%%%
\section{Pure spinor strings}
\frame{
\frametitle{Pure spinor strings}
In ten dimensions we can introduce spinor harmonics to separate constraints 
\begin{eqnarray}
\left\{d_\alpha , d_\beta\right\} = -2i (\gamma^m)_{\alpha \beta}\Pi_m &\to& 
	\left\{ \lambda^\alpha d_\alpha , \lambda^\beta d_\beta\right\} = -2i \, (\lambda \gamma^m \lambda)\,  \Pi_m
\nonumber
\end{eqnarray}
%\pause
%\begin{eqnarray}
%\left\{ \lambda^\alpha d_\alpha , \lambda^\beta d_\beta\right\} = 0 &\Leftrightarrow& \lambda \gamma^m \lambda = 0
%\nonumber
%\end{eqnarray}
Can be chosen to vanish ~~~ $\Leftrightarrow$ ~~~ $\lambda \gamma^m \lambda = 0$
%\pause

Fierz identity
\begin{eqnarray}
\lambda^\alpha \lambda^\beta = \sum_{p\mathrm{~odd}} (\gamma_{[p]})^{\alpha \beta} ( \lambda\gamma^{[p]}\lambda)= (\gamma_{m_1\dots m_5})^{\alpha \beta} (\lambda \gamma^{m_1\dots m_5}\lambda)
\nonumber
\end{eqnarray}

%\pause
defines self-dual 5-plane 

%\pause
$\Leftrightarrow$ Lagrangian sub-bundle of $T^\mathbb C M$ 

%\pause

$\Leftrightarrow$ ``pure spinor''
}

\frame{
\frametitle{Pure spinor strings}
Berkovits operator 
\begin{eqnarray}
Q=\oint \lambda^\alpha d_\alpha &\Rightarrow& Q^2=0
\nonumber
\end{eqnarray}
%\pause 
Cohomology at ghost number 1
%\pause
\begin{eqnarray}
U=\lambda^\alpha A_\alpha(x,\theta) \in H^1_Q &\Rightarrow& U = \Gamma^m A_m +\Gamma_\alpha \chi^\alpha - Q(G)
\nonumber
\end{eqnarray}
where
$
\Gamma^m =\lambda\gamma^m \theta$ and $\Gamma_\alpha =-\frac12 (\theta\gamma^{mpq}\theta) (\lambda \gamma_{mnq})_\alpha$ span the 0-momentum cohomology at ghost number 1.

%\pause
Cohomology corresponds to a photon $A_m$ and photino $\chi^\alpha$.
}

\frame{
\frametitle{Pure spinor strings}
Applications 
%\pause
\begin{itemize}
\item Compactification and relation to hybrid
%\pause
\item Holographic backgrounds
\item \dots
\end{itemize}
%\pause
Open problems
%\pause
\begin{itemize}
\item Looks like gauge fixing of a superstring with some fermionic symmetry. {\em Which string ??}
%\pause
\item Spectrum
\item \dots
\end{itemize}

}


%%%%%%%%%%%%%%%%%%%%%%%%%%%%%%%%%%%%%%%%%%%%%%%%%
\section{Conclusion}
\frame{
\frametitle{In other news...}
{\bf Hybrid}
\begin{itemize}
\item Compactification on $K3$, $K3\times T^2$ \cite{Linch:2006sh}, $CY_3$, $CY_4$, ...
\item Flux compactification on $CY$ \cite{Linch:2006ig,Linch:2008rw}, superpotentials, non-commutativity, ...
\item Superstring field theory, black holes, amplitudes, \dots
\end{itemize}

{\bf Pure spinor}
\begin{itemize}
\item Compactification on $CY$ \cite{Chandia:2009it}
\item $AdS_5 \times S^5$ background and possible proof of Maldacena duality
\item Integrability in the $AdS_5\times S^5$ sigma model \cite{Linch:2008nt}
\end{itemize}

%
%\begin{itemize}
%\item<1-> Superpotentials, non-commutativity, {\it et cetera}
%\item<2->
%\end{itemize}


}

\frame{
\frametitle{Fin}
\begin{centering}
{\Huge Thank you!}
\end{centering}
}

%%%%%%%%%%%%%%%%%%%%%%%%%%%%%%%%%%%%%%%%%%%%%%%
\begin{thebibliography}{99}

%\bibitem{AP13}
%\cite{Linch:2006ig}
\bibitem{Linch:2006ig}
  Wm~D.~Linch, III and B.~C.~Vallilo,
  ``Hybrid formalism, supersymmetry reduction, and Ramond-Ramond fluxes,''
  JHEP {\bf 0701}, 099 (2007)
  [arXiv:hep-th/0607122].
  %%CITATION = JHEPA,0701,099;%%
  
%%\cite{Berkovits:1999du}
%\bibitem{Berkovits:1999du}
%  N.~Berkovits,
%  ``Quantization of the type II superstring in a curved six-dimensional
%  background,''
%  Nucl.\ Phys.\  B {\bf 565}, 333 (2000)
%  [arXiv:hep-th/9908041].
%  %%CITATION = NUPHA,B565,333;%%

%\cite{Linch:2006sh}
\bibitem{Linch:2006sh}
  Wm~D.~Linch, III and B.~C.~Vallilo,
  ``Covariant N = 2 heterotic string in four dimensions,''
  JHEP {\bf 0703}, 082 (2007)
  [arXiv:hep-th/0611105].
  %%CITATION = JHEPA,0703,082;%%

%%\cite{Berkovits:1994wr}
%\bibitem{Berkovits:1994wr}
%  N.~Berkovits,
%  ``Covariant quantization of the Green-Schwarz superstring in a Calabi-Yau
%  background,''
%  Nucl.\ Phys.\  B {\bf 431}, 258 (1994)
%  [arXiv:hep-th/9404162].
%  %%CITATION = NUPHA,B431,258;%%
  
%\cite{Linch:2008rw}
\bibitem{Linch:2008rw}
  Wm~D.~Linch, III, J.~McOrist and B.~C.~Vallilo,
  ``Type IIB Flux Vacua from the String Worldsheet,''
  JHEP {\bf 0809}, 042 (2008)
  [arXiv:0804.0613 [hep-th]].
  %%CITATION = JHEPA,0809,042;%% 

%%\cite{Berkovits:2001tg}
%\bibitem{Berkovits:2001tg}
%  N.~Berkovits, S.~Gukov and B.~C.~Vallilo,
%  ``Superstrings in 2D backgrounds with R-R flux and new extremal black
%  holes,''
%  Nucl.\ Phys.\  B {\bf 614}, 195 (2001)
%  [arXiv:hep-th/0107140].
%  %%CITATION = NUPHA,B614,195;%%

%\cite{Chandia:2009it}
\bibitem{Chandia:2009it}
  O.~Chandia, Wm~D.~Linch, III and B.~C.~Vallilo,
  ``Compactification of the Heterotic Pure Spinor Superstring I,''
  JHEP {\bf 0910}, 060 (2009)
  [arXiv:0907.2247 [hep-th]].
  %%CITATION = JHEPA,0910,060;%%

%\cite{Linch:2008nt}
\bibitem{Linch:2008nt}
  Wm~D.~Linch, III and B.~C.~Vallilo,
  ``Integrability of the Gauged Linear Sigma Model for $AdS_5\times S^5$''
  JHEP {\bf 0911}, 007 (2009)
  [arXiv:0804.4507 [hep-th]].
  %%CITATION = JHEPA,0911,007;%%


\end{thebibliography}


\end{document}